\section{Бюджетное открытие и компиляция препятствий}\label{sec:obstruction}
Конечный эталонный список $\CanonRef$ определяет, что должен восстановить поиск; его предварительное вычисление противоречило бы цели бюджетного поиска. Используем оплачиваемый исполнитель, корректность которого определяется относительно $\CanonRef$ без вызова эталонного перечисления на каждом шаге.

\subsection{Оплаченные префиксы и сертифицированное отсечение}
Детерминированная функция котировки сопоставляет состоянию $u$ отсутствие шага или пару $(\eta, k)$ с положительной натуральной платой $\eta$ и отложенным вычислением следующего состояния $k:\mathrm{Unit}\to U$. Счёт $(u,B,D)$ хранит остаток средств $B$ и зарегистрированные расходы $D$. Если $\eta\le B$, исполнитель вызывает $k$ и меняет счёт на $(k(),B-\eta,D+\eta)$; иначе останавливается, не вызывая $k$. Ограниченный запуск считает шаги отдельно от средств; запаса в $B$ шагов достаточно, чтобы достичь состояния без доступного по цене следующего шага. Индукция даёт точное сохранение $B'+D'=B+D$ и любого индуктивного инварианта котируемых переходов (\rc{accounting}, \rc{invariant}, \rc{budget\_bound} в \rc{BudgetedDiscovery.lean}).

\begin{theorem}[Точное бюджетное продолжение]\label{thm:resume}
Для фиксированной детерминированной функции котировки, не зависящей от бюджета и имеющей положительные платы, пусть $\mathsf E$ --- исполнитель, а $\mathsf U_b$ прибавляет $b$ к остатку средств. Тогда
\[
 \mathsf E\mathsf U_{b_1+b_2}(s)=\mathsf E\mathsf U_{b_2}(\mathsf E\mathsf U_{b_1}(s)).
\]
Это равенство включает состояние, остаток и расходы даже при циклической системе переходов.
\end{theorem}
\begin{proof}
Каждый префикс, доступный при меньшем бюджете, доступен и при большем; оплата коммутирует с добавлением разности. Применим индукцию по оплаченному префиксу, считая случаи недоступного и отсутствующего шага стационарными. Положительные платы ограничивают число оплаченных шагов остатком средств, откуда следует тождество для исполнителя до исчерпания доступных шагов. Это \rc{execute\_topUp\_add} в \rc{ResumeLaws.lean}. Теорема не влечёт семантического исчерпания поиска: цикл может просто израсходовать бюджет.
\end{proof}
В поиске морфизмов частичное назначение портов $\rho$ задаёт область продолжающих его морфизмов. Сертифицированное правило отсечения содержит булев тест $k(\rho)$, плату и свидетельство
\[
 \forall M,\rho,\quad \operatorname{Extends}(M,\rho)\Longrightarrow k(\rho)=\mathrm{true}.
\]
Ложный результат исключает всю область. Правила могут проверять инъективность, цены портов, глубину или совместимость чтений; корректные правила компонуются. Поиск с отсечением использует тот же конструктор канонической пары, что и раздел~3.

\begin{theorem}[Отсечение сохраняет каноническое членство]\label{thm:prune}
Для каждого сертифицированного правила и морфизма $x$ членство в полученном исчерпывающем списке с отсечением эквивалентно $x\in\CanonRef$.
\end{theorem}
\begin{proof}
Сохранённое полное назначение проверяется эталонным конструктором пары, что даёт корректность. Для полноты проследим префиксы карты портов любого канонического представителя; корректность каждого правила обеспечивает сохранение всех этих префиксов. Затем рассматривается эталонная карта действий и получается тот же канонический сертификат. Это \rc{prunedDiscover\_iff} в \rc{CertifiedPruning.lean}; равенство порядка и кратности списков не требуется.
\end{proof}
Корректное отсечение не обязано экономить ресурсы: при $n_1\le n_0$ посещённых узлах, единой цене узла $c$ и цене правила на узел $k$ оно окупается в точности при $n_1 k<(n_0-n_1)c$ (\rc{pruning\_profitable\_iff}). Корректное правило без эффекта повышает стоимость, если за него взимается плата. Более развитый исполнитель ниже отдельно учитывает фактические вызовы правил с досрочным прекращением проверки.

\subsection{Финансируемый фронтир поиска}
Фронтир $f=(J,A,\Lambda)$ хранит оставшиеся задания $J$, выданные канонические сертификаты $A$ и оплаченные поисковые квитанции $\Lambda$. Задание либо расширяет область назначений портов с оставшейся очередью правил, либо обрабатывает оставшийся суффикс карт действий. Корректный фронтир содержит только канонические сертификаты и покрывает каждый ещё не выданный канонический сертификат хотя бы одним оставшимся заданием; точный предикат дан в разделе~6.1. Разбиение задания, исключение сертифицированной области или проверка пары сохраняют этот предикат и дописывают новый сертификат, если он найден. Доказательство использует локальное покрытие и корректность, а не исполнение эталонного перечисления.

Пусть $P(C,f)$ --- натуральный потенциал оставшихся обязательств при фиксированной библиотеке правил и базовых тарифах. Бюджеты его заданий ограничивают полную работу ветвления, включая платы за правила. Каждый котируемый платный шаг удовлетворяет
\begin{equation}\label{eq:funded-potential}
 P(C,f')+\eta\le P(C,f),\qquad\eta>0.
\end{equation}
\begin{theorem}[Финансируемое исчерпание]\label{thm:funded-search}
Начиная с корректного фронтира при $B\ge P(C,f)$, исполнитель исчерпывает задания, сохраняет прежние сертификаты как префикс списка и возвращает список с членством в точности как в $\CanonRef$. Пустой список сертификатов при исчерпанном фронтире доказывает отсутствие морфизма.
\end{theorem}
\begin{proof}
Уравнение~\eqref{eq:funded-potential} сохраняет финансирование после платежа и строго уменьшает потенциал. При остановке наличие оставшегося задания означало бы доступный по цене положительный шаг, что противоречит остановке. Локальная корректность индуктивна, а пустой список заданий устраняет дизъюнкт оставшегося покрытия. Сохранение префикса следует по индукции из дописывания сертификатов. Это \rc{quote\_potential}, \rc{execute\_funded\_exhausts}, \rc{execute\_full}, \rc{execute\_certificate\_prefix} и \rc{exhausted\_impossible} в \rc{BudgetedPruning.lean}.
\end{proof}
Без предпосылки финансирования классификатор различает найденный свидетель, пустой исчерпанный фронтир с доказательством невозможности и неопределённый исход при пустом списке сертификатов и неисчерпанном фронтире. Возвращённый счёт допускает продолжение. Найденный до исчерпания свидетель не означает полноту перечисления.

\subsection{От пустой строки к переиспользуемому правилу}
Для целевого мира $w$ частичное назначение портов $\rho$ определяет строку совместимых чтений
\[
 Y_w(\rho)=\{y\in X_S: \forall(q,q')\in\rho,\ r_S(q,y)=r_T(q',w)\}.
\]
Если $Y_w(\rho)=\varnothing$, добавление ограничений не может восстановить элемент. Компилятор препятствий удаляет ограничения, сохраняя пустоту строки, и получает подмножество $K\subseteq\rho$, которое остаётся несовместимым, но перестаёт быть таковым при удалении любого из значений ограничений. Таким образом, $K$ неприводимо по включению, но не обязательно имеет минимальную мощность.
\begin{theorem}[Сертифицированная компиляция препятствий]\label{thm:core}
Для пустой строки чтений \rc{extract} строит такое неприводимое $K$ со свидетельством. Ни одно назначение, содержащее $K$, не допускает продолжающего морфизма. Скомпилированный тест отклоняет в точности назначения, содержащие $K$, и удовлетворяет контракту сертифицированного правила.
\end{theorem}
\begin{proof}
Процедура сокращения просматривает исходные ограничения, удаляя значение, если несовместимость сохраняется. Если удаление не удалось, монотонность несовместимости гарантирует, что последующие удаления не могут снова сделать результат этого удаления несовместимым. Поэтому оставшееся подмножество неприводимо. Если бы существовал морфизм, продолжающий назначение, содержащее $K$, его карта миров в $w$ дала бы элемент пустой строки. Булев тест включения точен по отражению членства. См. \rc{shrink\_spec}, \rc{extract}, \rc{core\_excludes}, \rc{compileCore\_rejects} в \rc{ImpasseCompilation.lean}.
\end{proof}
Перенос ядра на новую пару задач допускается при наличии карты ограничений, отображения новых исходных миров в старые и доказательства отражения чтений для каждого перенесённого ограничения (\rc{transportCore}); неприводимость и экономическое превосходство перенесённого правила не гарантируются.

Обучающийся исполнитель делает компиляцию эндогенной. Состояние поиска включает активную библиотеку ядер, исходные и сокращённые события обучения, ожидающую компиляцию и оплаченные квитанции. Поиск в библиотеке, зондирование, компиляция и базовое продвижение поиска --- отдельные переходы с положительными платами. Ядро вставляется с поглощением: если старое ядро уже исключает всю его область, новое не устанавливается; иначе удаляются ядра, ограничения которых содержат новое ядро. Итоговая исключаемая область в точности равна объединению прежней области и области нового ядра (\rc{insertCore\_rejects}); предыдущие события обучения сохраняются в истории. Теоремы \rc{execute\_correct}, \rc{execute\_trace}, \rc{execute\_split} и \rc{execute\_exhausted\_iff} в \rc{LearningDiscovery.lean} сохраняют семантические и учётные гарантии. Исчерпание условно и не гарантируется произвольным лимитом обучения.

Эта конструкция делает отрицательные результаты операционно переиспользуемыми. Она не устанавливает универсально более дешёвого обучающегося исполнителя. В объявленном конечном стенде обучение было дороже фиксированного контроллера правила чтения во всех 515 полностью завершённых случаях. Поэтому следующая задача --- выбор самого исследовательского механизма, а не безусловное обучение.
